% Generated draft: requires mathematical and source review.
\documentclass[11pt]{article}
\usepackage[margin=25mm]{geometry}
\usepackage{fontspec}
\setmainfont{Latin Modern Roman}
\usepackage{amsmath,amssymb,mathtools,unicode-math}
\setmathfont{Latin Modern Math}
\usepackage{enumitem,needspace,xurl,hyperref}
\setlength{\parindent}{0pt}
\setlength{\parskip}{5pt}
\setlength{\emergencystretch}{3em}
\newcommand{\N}{\mathbb{N}}
\newcommand{\Z}{\mathbb{Z}}
\newcommand{\Q}{\mathbb{Q}}
\newcommand{\R}{\mathbb{R}}
\newcommand{\C}{\mathbb{C}}
\newcommand{\F}{\mathbb{F}}
\newcommand{\A}{\mathbb{A}}
\newcommand{\PP}{\mathbb{P}}
\newcommand{\E}{\mathbb{E}}
\newcommand{\eps}{\varepsilon}
\newcommand{\dd}{\,\mathrm{d}}
\begin{document}
\section*{UnsolvedMath ID: 701456}
\textbf{Source problem code:} catalogue entry
\subsection{Definitions and mathematical statement}
For this draft, interpret a formation as a nonempty class \(\mathcal F\) of finite groups, invariant under isomorphism and closed under homomorphic images and finite subdirect products. Here a subdirect product of \(H_1,\ldots,H_n\), for \(n\geq 1\), is a subgroup \(H\leq H_1\times\cdots\times H_n\) whose projection onto each factor is surjective. Closure under finite subdirect products means that every such \(H\) belongs to \(\mathcal F\) whenever all the factors belong to \(\mathcal F\).

For a finite group \(G\), define its generated formation by
\[
\operatorname{form}(G)=\bigcap\{\mathcal F:\mathcal F\text{ is a formation and }G\in\mathcal F\}.
\]
A subformation of \(\operatorname{form}(G)\) is a formation \(\mathcal H\) satisfying \(\mathcal H\subseteq\operatorname{form}(G)\). Write
\[
\mathscr S(G)=\{\mathcal H:\mathcal H\text{ is a formation and }\mathcal H\subseteq\operatorname{form}(G)\}.
\]
The problem asks whether there exists a single algorithm which, for every finite group \(G\) given by its multiplication table, halts and correctly decides whether \(\mathscr S(G)\) is finite.

The supplied record does not define formation, subformation, or \(\operatorname{form}(G)\). The finite-group formation convention above is an explicit interpretation requiring confirmation against the source; the multiplication-table presentation makes the algorithmic input precise.
\subsection{Short English statement}
V. P. Burichenko asks whether an algorithm can decide, for every finite group \(G\), whether the formation generated by \(G\) contains only finitely many subformations. This draft uses the stated finite-group formation convention, which requires source confirmation.
\subsection{Sources}
\begin{itemize}
\item Local UnsolvedMath registry, permanent ID \(701456\), exact source problem code \texttt{catalogue entry}, title \emph{Kourovka Notebook Problem 18.27}; supplied upstream revision \texttt{d371ee917ac9647fbb06cd6fc529deec6f3dda82}. The registry labels the problem open; this is supplied status information, not an independently verified assessment.
\item Supplied source attribution: \emph{The Kourovka Notebook (Unsolved Problems in Group Theory)}, 18th issue (2014), problem 18.27, proposed by V. P. Burichenko. The background also identifies the 21st issue edition (2026), arXiv:\texttt{1401.0300v48}.
\item The supplied collection description attributes editing of the Kourovka Notebook, Issues 1--21 (1965--2026), to E. I. Khukhro and V. D. Mazurov. These bibliographic details and the cited source have not been independently checked for this draft.
\end{itemize}
\end{document}
